% !TeX program = pdflatex
% papers/conversa_rede_dual.tex — Rede neural dual no corpus (leve; depois histerese).
\documentclass[11pt,a4paper]{article}
\usepackage[utf8]{inputenc}\usepackage[T1]{fontenc}\usepackage[brazil]{babel}
\usepackage[margin=2.4cm]{geometry}
\usepackage{xcolor}
\usepackage[hidelinks]{hyperref}
\newcommand{\code}[1]{\texttt{\small #1}}
% ─────────────────────────────────────────────────────────────────────────────
%  papers/gkcapa.tex — a capa do Reino, mínima, para os papers em `article`.
%
%  O estilo.tex completo é do `report` (títulos de capítulo, skak, …). Aqui fica
%  só o que a capa precisa, para o paper continuar a compilar sozinho com
%  pdflatex. O tradutor próprio (tests/tex) carrega o estilo.tex da raiz / o
%  symlink papers/estilo.tex e avalia o mesmo `\gkcapa`.
% ─────────────────────────────────────────────────────────────────────────────
\definecolor{ouro}{HTML}{B8912F}
\definecolor{ouroclaro}{HTML}{F3E9CE}
\definecolor{tinta}{HTML}{1E1B16}
\definecolor{regua}{HTML}{6E6858}
\providecommand{\gknota}{\fontsize{7.62}{11.05}\selectfont}
\providecommand{\gktexto}{\fontsize{10.50}{15.22}\selectfont}
\providecommand{\gksub}{\fontsize{12.33}{17.87}\selectfont}
\providecommand{\gksec}{\fontsize{14.47}{20.98}\selectfont}
\providecommand{\gkcap}{\fontsize{16.99}{24.63}\selectfont}
\providecommand{\gktit}{\fontsize{23.42}{33.95}\selectfont}
\providecommand{\Prj}{\mathbb{P}}
\providecommand{\Trf}{\mathcal{T}}
\providecommand{\gkcapa}[3]{%
\title{\vspace{-0.6cm}
{\color{ouro}\rule{\textwidth}{1.4pt}}\\[6mm]
\textbf{\gktit\color{tinta} Reino Dourado}\\[3mm]{\gksec\itshape\color{regua} Golden Kingdom}\\[8mm]
{\gkcap\scshape\color{ouro} #1}\\[5mm]
{\gksub\color{regua} #2}\\[2mm]
{\gktexto\color{regua}\itshape #3}\\[7mm]
{\color{ouro}\rule{\textwidth}{0.6pt}}\\[9mm]{\gknota\color{regua}\begin{minipage}{\textwidth}\setlength{\parindent}{0pt}%
\textbf{Aviso de Isenção Algorítmica:} O universo, as leis e as entidades contidas nestas páginas ---
incluindo felinos matriciais, aves de rapina algébricas e amuletos de controle de fluxo --- são
construções de pura ficção formal. Esta obra não descreve a realidade física, não serve como
engenharia reversa do cosmos e não constitui doutrina religiosa. Qualquer semelhança com bugs reais do seu
sistema operacional é mera coincidência (ou falta de reza). Prossiga por sua própria conta e
risco.\end{minipage}}}
\author{}\date{}}

\begin{document}
\gkcapa{Rede neural dual}
       {P $\cdot$ $\mathcal{H}$ $\cdot$ H na interface}
       {estaca, banda, memória --- conjugação, não atrator}
\maketitle
\begin{abstract}
Pares curtos do Teorema da Rede Dual.
Detalhe: \code{papers/corpo\_peano.tex} \S rede-dual.
UI: \code{app/src/rede\_dual.js}; medidor \code{tests/rede\_dual.js}.
\end{abstract}

\section{o que e a rede dual}
$P\cdot\mathcal{H}\cdot H$: estaca $W=-I$, banda de histerese, Hopfield $=$ memória $Y$.
$\lambda^{-}$ vem da conjugação reversível, não do atrator.
\code{papers/corpo\_peano.tex} \S rede-dual.

\section{o que e a rede neural dual}
Teorema: retenção na banda; $F_H=D\circ F_P^{-1}\circ D^{-1}$; $\lambda_P+\lambda_H=0$.
Hopfield $=$ memória da volta, não inversão. Pesos $W=-I$ pela Lei~1.
\code{papers/corpo\_peano.tex} \S rede-dual.

\section{o que e a conjugacao reversivel}
$F_H=D\circ F_P^{-1}\circ D^{-1}$ com $D(x)=2c-x$ (estaca Dual Sort).
Hopfield $\neq$ inversão: atracção $\neq$ reversibilidade.
\code{papers/corpo\_peano.tex} \S rede-dual; \code{tests/rede\_dual.js} §R5--§R7.

\section{o que e o estado hibrido}
$X=(x,h)$; $\mathcal{F}\colon(x,h)\mapsto(x',h')$.
Na banda: $(x',h')=(x,h)$. Fora: frente $F_P$ ou volta $F_H$.
\code{papers/corpo\_peano.tex} \S rede-dual.

\section{o que e retencao neural}
$|u|\le\Delta\Rightarrow h'=h$; sequência na banda conserva $h_n=h_0$.
\code{papers/corpo\_peano.tex} thm:retencao-neural.

\section{hopfield e a inversa?}
Não. Hopfield realiza a memória da volta; a dualidade reversível garante $\lambda^{-}$.
\code{papers/corpo\_peano.tex} \S rede-dual.

\section{o que mede o metronomo na rede}
$\lambda^{+}+\lambda^{-}=0$ quando a volta fecha --- atestação, não oráculo.
\code{papers/corpo\_peano.tex}; UI \code{app/src/assistente.js}.

\section{mostra a rede dual}
\code{papers/corpo\_peano.tex} \S rede-dual --- depois histerese e leitura complexa.
Medidor: \code{tests/rede\_dual.js}.

\section{falemos da rede dual}
Três operadores distintos: percepção, borda, memória.
\code{papers/corpo\_peano.tex} \S rede-dual.

\section{a conjugacao fechou no experimento?}
Sim, no modelo afim Dual Sort: $DF_H\,DF_P=I$, $F_H\circ F_P=\mathrm{id}$,
$\lambda_P+\lambda_H=0$ em dim.~1 e~2 (\code{tests/rede\_dual.js} §R5--§R7;
Prop.\ fecho experimental em \code{papers/corpo\_peano.tex}).
Hopfield continua memória $Y$, não a prova de inversão.

\section{a rede dual e o teorema central?}
Não é o Teorema Central: é o controlo de histerese sob Maestro/Metrónomo.
Cadeia: Central $\to$ Maestro/Metrónomo $\to$ controlo $\to$ rede dual.
\code{papers/corpo\_peano.tex} Prop.\ leitura operacional.

\section{o que e o teorema central no peano}
Hurwitz (contar) $\leftrightarrow$ Gentil (integrar) pela soma reversível ---
estrela $\nu\circ\nu=\mathrm{id}$. Citado do Corpo estelar; na torre vira
$\pi_k$, Maestro/Metrónomo, depois controlo.
\code{papers/corpo\_peano.tex}; \code{teoria.tex} \S central.

\section{como o teorema central sobe a torre}
Estrutura $\pi_k$ $\to$ dinâmica/métrica (Maestro projecta; Metrónomo
atesta $\lambda^{+}+\lambda^{-}=0$) $\to$ controlo $(P,\mathcal{H},H)$
$\to$ Conservatório/Pera/Batuta.
Mesma forma da estrela; camadas distintas.
\code{papers/corpo\_peano.tex} Prop.\ central-operacional.

\section{o que e o conservatorio no peano}
Andar acima: estrutura superior conservada (norma Gentil, cone).
Não é a orquestra --- é o depósito que a projecção não inventa.
\code{papers/corpo\_peano.tex} \S musica.

\section{o que e a pera na batuta}
Base da batuta: cone acoplado em $s=0$; encaixe
Conservatório$\hookrightarrow$Batuta. Não é $I$.
\code{papers/corpo\_peano.tex}.

\section{como a rede dual entra na musica}
A batuta canaliza $I$; a rede dual escolhe $F_P$/retenção/$F_H$;
o Metrónomo atesta $\lambda^{+}+\lambda^{-}=0$ relativamente a $\Pi$.
\code{papers/corpo\_peano.tex} Prop.\ rede-no-canal.

\section{o que e a partitura pi}
$\Pi=$ assinatura $+$ notação: o \emph{quê} do Maestro.
Música $=$ realização; partitura $=$ transporte. Não funda leis.
\code{papers/corpo\_peano.tex} Def.\ partitura; \code{papers/partitura.tex}.

\section{como a partitura entra na assistente}
\code{partituraFala} extrai $\Pi$; batuta dá $I$; rede dual escolhe o ramo;
Metrónomo lê $r$ sob $\Pi$.
\code{app/src/maestro.js}; Prop.\ pi-cadeia em \code{papers/corpo\_peano.tex}.

\section{mostra a partitura}
\code{papers/partitura.tex} --- realização gráfica; produto
\code{partitura-wasm.pdf}.

\end{document}
